\section{Recovery of the state and its dynamics}
\label{nuclear:3}
\subsection{Recoverable representation of transport}

Let \(C\) be the unitary realization of memory transport on \(\mathcal H\). The nine-phase chart, its uniform inclusion, and its compression are

\[
S_C(v_0,\ldots,v_8)=(Cv_8,v_0,\ldots,v_7),\qquad
Jv=(v,\ldots,v)/3,\qquad T=(8I+C)/9.
\]

Let \(\eta=(I-JJ^*)S_CJ\). Its action is

\[
\eta v=\frac1{27}\bigl(8(C-I)v,-(C-I)v,\ldots,-(C-I)v\bigr).
\]

Normalization of nine copies determines \(1/3\); the distribution of the seam determines \(8/81\) in

\begin{equation}
I-T^*T=\eta^*\eta=\frac8{81}(I-C)^*(I-C).
\label{nuc03:eq:1}
\end{equation}

Write \(P=\operatorname{proj}\ker(I-C)\). The preceding theorem proves strong convergence \(T^N\to P\). Therefore,

\begin{equation}
\mathcal Vv=(Pv,\eta v,\eta Tv,\eta T^2v,\ldots)
\quad\text{in}\quad
\mathcal R=P\mathcal H\oplus\ell^2(\mathbb N_0;\mathcal H^9)
\label{nuc03:eq:2}
\end{equation}

is an isometry. Denote its closed image by \(\mathcal M\).

On the record space, define the block shift

\begin{equation}
\mathcal B(a,m_0,m_1,m_2,\ldots)=(a,m_1,m_2,\ldots).
\label{nuc03:eq:3}
\end{equation}

The first component preserves the fixed sector; the remaining components are the complete vector records, including their correlations.

\begin{theorem}
\label{nuc03:transporte}
The original transport has the exact representation

\begin{equation}
\boxed{\mathcal V C=(9\mathcal B-8I)\mathcal V.}
\label{nuc03:eq:4}
\end{equation}

The restriction \(\mathcal C=(9\mathcal B-8I)|_{\mathcal M}\) is unitary on \(\mathcal M\), and

\begin{equation}
C=\mathcal V^*\mathcal C\mathcal V.
\label{nuc03:eq:5}
\end{equation}
\end{theorem}

\begin{proof}
Since \(PT=P\), definition~\eqref{nuc03:eq:2} gives, componentwise,

\[
\mathcal VT v=(Pv,\eta Tv,\eta T^2v,\ldots)=\mathcal B\mathcal Vv.
\]

The algebraic identity \(C=9T-8I\) implies \eqref{nuc03:eq:4}. The isometry \(\mathcal V:\mathcal H\to\mathcal M\) is surjective by definition of its image. Equality \eqref{nuc03:eq:4} identifies the restriction with \(\mathcal V C\mathcal V^{-1}\), which is unitary because \(C\) is unitary. This also proves \(\mathcal C\mathcal M=\mathcal M\) and \eqref{nuc03:eq:5}.
\end{proof}

Restriction to the subspace of compatible records is essential: \(9\mathcal B-8I\) on the whole of \(\mathcal R\) is not unitary. The theorem characterizes transport on records effectively arising from the same linear trajectory of the documented chart.

\begin{corollary}
For \(v,w\in\mathcal H\) and \(k\in\mathbb Z\),

\begin{equation}
\langle C^kv,w\rangle
=\langle\mathcal C^k\mathcal Vv,\mathcal Vw\rangle.
\label{nuc03:eq:6}
\end{equation}

All cross moments of the realized dynamics are thus preserved. Negative powers belong to the inverse of \(\mathcal C\) on \(\mathcal M\), rather than to an inversion of the unilateral shift on the ambient space.
\end{corollary}

\begin{corollary}
For every bounded observable \(A\), its realization on the image is \(\widehat A=\mathcal V A\mathcal V^*|_{\mathcal M}\). Products, adjoints, and commutators are preserved. In particular,

\begin{equation}
[A,C]=0\quad\Longleftrightarrow\quad[\widehat A,\mathcal C]=0.
\label{nuc03:eq:7}
\end{equation}
\end{corollary}

The proof uses \(\mathcal V^*\mathcal V=I\) and the fact that \(\mathcal V\mathcal V^*\) is the identity on \(\mathcal M\). Symmetry conservation is expressed through complete transport of its operators, in addition to an equality of norms.

\subsection{Explicit reconstruction and record depth}

Let \(m_j=\eta T^jv\). The telescopic identity gives, for each \(N\),

\begin{equation}
\boxed{v=T^{*N}T^Nv+\sum_{j=0}^{N-1}T^{*j}\eta^*m_j.}
\label{nuc03:eq:8}
\end{equation}

Passing to the strong limit,

\begin{equation}
\boxed{v=Pv+\sum_{j=0}^{\infty}T^{*j}\eta^*m_j.}
\label{nuc03:eq:9}
\end{equation}

Partial reconstruction with fixed terminal component,
\(v_N=Pv+\sum_{j<N}T^{*j}\eta^*m_j\), has the exact error

\begin{equation}
v-v_N=T^{*N}T^N(I-P)v.
\label{nuc03:eq:10}
\end{equation}

The series converges for every state. If the complete record has an error \(\delta r\) measured in the norm of \(\mathcal R\), its reconstruction through \(\mathcal V^*\) changes by at most \(\|\delta r\|\) in norm, because \(\|\mathcal V^*\|=1\). This is uniform stability of the complete readout.

\subsubsection{Difference between finite and unlimited depth}

For the bilateral shift \(Ce_k=e_{k+1}\) on \(\ell^2(\mathbb Z)\), one has \(P=0\). Consider

\[
v_L=L^{-1/2}\sum_{k=1}^{L}e_k.
\]

Then \(\|(I-C)v_L\|^2=2/L\). Commutation of \(T\) with \(I-C\), contractivity of \(T\), and \eqref{nuc03:eq:1} imply

\begin{equation}
\sum_{j=0}^{N-1}\|\eta T^jv_L\|^2
\leq\frac{16N}{81L}.
\label{nuc03:eq:11}
\end{equation}

For any fixed depth \(N\), the right-hand side tends to zero as \(L\to\infty\), although \(\|v_L\|=1\). The fixed-depth memory record, separated from its terminal component, has no uniform lower bound. By contrast, the unlimited record has norm exactly one for each \(v_L\).

The difference corresponds to two distinct quantifiers: every state is preserved upon completion of its record; a prescribed depth does not uniformly recover every state without its terminal component. Conservation at arbitrary depth here has an exact functional expression.

\subsubsection{Vector information and intensity}

The theorem preserves \(m_j\), not merely \(\|m_j\|^2\). For example, in the bilateral realization, \(v=e_0\) and \(w=e_1\) have identical intensities at every depth because the shift commutes with \(T\) and transports \(\eta\) unitarily. Their vector records differ, and \eqref{nuc03:eq:9} recovers the two distinct states. Intensities serve the balance; records with their correlations additionally serve reconstruction.

\subsection{Naturality of recovery under refinement}

Let \(\iota_n:\mathcal H_n\to\mathcal H_{n+1}\) be the refinement isometries obtained from the common emissions, and let \(C_n\) be unitary realizations with

\begin{equation}
C_{n+1}\iota_n=\iota_nC_n.
\label{nuc03:eq:12}
\end{equation}

The relation also holds for adjoints: multiplication by \(C_{n+1}^*\) and \(C_n^*\) gives \(C_{n+1}^*\iota_n=\iota_nC_n^*\). The images are therefore reducing. The polynomial formulas for \(T_n\), \(\eta_n\), and the strong limit of \(T_n^N\) give

\begin{equation}
T_{n+1}\iota_n=\iota_nT_n,\quad
\eta_{n+1}\iota_n=\iota_n^{\oplus9}\eta_n,\quad
P_{n+1}\iota_n=\iota_nP_n.
\label{nuc03:eq:13}
\end{equation}

Let \(\mathfrak I_n\) be the refinement applying \(\iota_n\) to the terminal component and \(\iota_n^{\oplus9}\) to every record. Then

\begin{equation}
\boxed{\mathcal V_{n+1}\iota_n=\mathfrak I_n\mathcal V_n,\qquad
\mathcal B_{n+1}\mathfrak I_n=\mathfrak I_n\mathcal B_n.}
\label{nuc03:eq:14}
\end{equation}

\begin{proof}
Each component of the first equality is an identity from \eqref{nuc03:eq:13}, iterated to the corresponding power of \(T_n\). The second equality states that refining every block commutes with shifting the block index, as verified directly. Both maps are bounded, so the equalities pass from finitely supported records to their completion.
\end{proof}

By \eqref{nuc03:eq:4} and \eqref{nuc03:eq:14}, reconstruction of the dynamics also commutes with refinement. The genealogical-depth index \(n\) and the successive-compression index \(j\) remain distinct. The equality expresses their compatibility when they act through the documented squares, rather than identifying them by name.

\subsection{Two exact regimes of memory distribution}

For a unit state \(v\) with \(Pv=0\), define

\begin{equation}
a_N=\|T^Nv\|^2,\qquad b_N=\|\eta T^Nv\|^2=a_N-a_{N+1}.
\label{nuc03:eq:15}
\end{equation}

One has \(a_0=1\), \(a_N\to0\), \(b_N\geq0\), and \(\sum b_N=1\). The weights \(b_N\) describe the distribution of the memory record across depths. This index counts compressions; identifying it with a physical duration would require application of the corresponding temporal reader.

\subsubsection{Exceptional transport of order three}

The fixed-point-free orthogonal operator \(g_c\) satisfies \(g_c^2+g_c+I=0\). For the explicit composition \(C=g_c\), the preceding development obtained

\begin{equation}
T^*T=\frac{19}{27}I,
\qquad a_N=\left(\frac{19}{27}\right)^N,
\qquad b_N=\frac8{27}\left(\frac{19}{27}\right)^N.
\label{nuc03:eq:16}
\end{equation}

The mean depth, counting the first record as depth one, is

\begin{equation}
\sum_{N\geq0}(N+1)b_N=\frac{27}{8}.
\label{nuc03:eq:17}
\end{equation}

This is a result of this composition of the nonadic chart with the already constructed exceptional transport; it preserves the domain and does not identify \(g_c\) with complete temporal holonomy merely from its order.

\subsubsection{Bilateral record of revolutions}

For \(C e_k=e_{k+1}\) and \(v=e_0\), the binomial formula determines every term:

\begin{equation}
T^Ne_0=9^{-N}\sum_{k=0}^{N}\binom{N}{k}8^{N-k}e_k,
\quad
a_N=81^{-N}\sum_{k=0}^{N}\binom{N}{k}^{\!2}64^{N-k}.
\label{nuc03:eq:18}
\end{equation}

This formula is exact at any depth. The subsequent circular representation gives

\begin{equation}
a_N=\frac1{2\pi}\int_{-\pi}^{\pi}
\left(\frac{65+16\cos\theta}{81}\right)^N\,d\theta.
\label{nuc03:eq:19}
\end{equation}

\begin{theorem}
In this bilateral realization,

\begin{equation}
\boxed{a_N\sim\frac9{4\sqrt{2\pi}}N^{-1/2},\qquad
b_N\sim\frac9{8\sqrt{2\pi}}N^{-3/2}.}
\label{nuc03:eq:20}
\end{equation}
\end{theorem}

\begin{proof}
Write \(r(\theta)=(65+16\cos\theta)/81\) and \(q=8/81\). Near zero,

\[
r(\theta)=1-q\theta^2+O(\theta^4),\qquad
\log r(\theta)=-q\theta^2+O(\theta^4).
\]

For a sufficiently small fixed neighbourhood, positive constants \(c_1,c_2\) bound \(-\log r(\theta)\) between \(c_1\theta^2\) and \(c_2\theta^2\). Outside this neighbourhood, \(r\leq\rho<1\). The substitution \(x=\sqrt N\theta\), the Gaussian bound, and dominated convergence give

\[
\lim_{N\to\infty}\sqrt N\,a_N
=\frac1{2\pi}\int_{\mathbb R}e^{-qx^2}\,dx
=\frac1{2\sqrt{\pi q}}=\frac9{4\sqrt{2\pi}}.
\]

For the second limit, apply the same procedure directly to

\[
b_N=\frac1{2\pi}\int_{-\pi}^{\pi}r(\theta)^N(1-r(\theta))\,d\theta.
\]

In the variable \(x\), \(N(1-r(x/\sqrt N))\to qx^2\), locally dominated by a constant times \(x^2\); outside the neighbourhood, the contribution is exponentially small even after multiplication by \(N^{3/2}\). Therefore,

\[
\lim_{N\to\infty}N^{3/2}b_N
=\frac q{2\pi}\int_{\mathbb R}x^2e^{-qx^2}\,dx
=\frac1{4\sqrt{\pi q}}=\frac9{8\sqrt{2\pi}}.
\]

This proof of the second limit uses its own integral, without differentiating an uncontrolled asymptotic equivalence.
\end{proof}

By Tonelli's theorem and the tail-sum identity for a positive sequence,

\begin{equation}
\sum_{N\geq0}(N+1)b_N=\sum_{N\geq0}a_N=+\infty.
\label{nuc03:eq:21}
\end{equation}

The entire norm is preserved in memory, but its mean depth diverges for this bilateral state. The distribution law depends on the effective transport representation: the exceptional sector of order three has a geometric rate, whereas the bilateral record considered here has an algebraic rate.

\subsection{Composition with Euler's identity and incidence}

The extended Euler identity of Section~\ref{euler-trit-generadores} realizes the three regimes through the Coxeter operator of \(A_2\), a nilpotent transport, and \(F_{\rm av}^2\), with their stated antecedents. Their normalized generators satisfy \(J_\tau^2=-\tau I\) and have zero trace. In their two-dimensional charts,

\[
C_\tau(t)=\exp(tJ_\tau),\qquad
\Omega=\begin{pmatrix}0&1\\-1&0\end{pmatrix},\qquad
C_\tau(t)^\mathsf T\Omega C_\tau(t)=\Omega.
\]

Zero trace gives \(\det C_\tau(t)=1\), and the two-dimensional identity \(M^\mathsf T\Omega M=(\det M)\Omega\) proves the final equality. For

\[
T_\tau=\frac{8I+C_\tau}{9},\qquad
D_\tau=\frac{\sqrt8}{9}(C_\tau-I),
\]

one obtains a common oriented conservation law:

\begin{equation}
\boxed{\Omega=T_\tau^\mathsf T\Omega T_\tau+D_\tau^\mathsf T\Omega D_\tau.}
\label{nuc03:eq:22}
\end{equation}

\begin{proof}
Expanding the two terms, the cross contributions of \(C_\tau\) have opposite coefficients. The remaining terms are \(72\Omega/81+9C_\tau^\mathsf T\Omega C_\tau/81=\Omega\).
\end{proof}

This alternating form corresponds to oriented area in the canonical plane. Its preservation permits comparison of the three regimes while respecting their differences. For the concrete representatives given by the Coxeter operator \(C\), \(I+N\), and \(F_{\rm av}^2\), respectively,

\begin{equation}
(\det T,\det D)=
\left(\frac{19}{27},\frac8{27}\right),\qquad
(1,0),\qquad
\left(\frac{89}{81},-\frac8{81}\right).
\label{nuc03:eq:23}
\end{equation}

The sum is one in all three cases. In the parabolic case, \(D\) has rank one and may preserve information although its area is zero. In the hyperbolic case, the complement has the opposite orientation; the balance concerns signed area. Theorem~\ref{nuc03:transporte} uses a unitary, positive realization: the hypotheses of each result are expressly retained.

The action and varying-form development proves the version between different charts, its integration over closed curves, and its compatibility with refinement. The incidence development composes reconstruction with \(F(k)=(\mu(k),\mathcal I P_0k/6)\), retaining all twelve coordinates of \(K\); it also transports operators and their local frames. The same conservation identity thus admits typed realizations of norm, incidence, and oriented action.
